\section{Cut measures and transport along a decomposition}
\label{sec:cuts}

The construction will prescribe a measure on the binary assignments of
each bag.  Our first task is to combine these measures into one global
cut metric.  Our second is to show that a cut measure present in every bag law
along a path can force separation in the glued metric, even when the
two vertices do not lie in a common bag.

Fix an integer $k\geq 2$.
Let $G=(V,E)$ be a finite connected simple undirected graph with a
nonempty vertex set and positive edge lengths $\ell:E\to(0,\infty)$.
Its shortest-path metric $d(u,v)$ is the minimum sum of edge lengths
along a path from $u$ to $v$, with $d(u,u)=0$.
For a nonempty set $A\subseteq V$, put
\[
 d(x,A)=\min_{a\in A}d(x,a),
 \qquad
 \Delta=\operatorname{diam}(V,d).
\]
A tree decomposition consists of a finite tree $T$ and bags
$B_t\subseteq V$, indexed by $t\in V(T)$, such that the bags cover $V$,
every graph edge has both endpoints in a bag, and
\[
 T_v=\{t\in V(T):v\in B_t\}
\]
induces a connected subtree for every $v\in V$.
The treewidth of $G$ is the minimum of $\max_t|B_t|-1$ over all its
tree decompositions. We assume $|B_t|\leq k$; thus the graph has
treewidth at most $k-1$.
Bags are regarded as indexed occurrences, even when two bags have the same
vertex set.

We may assume that all bags are nonempty.
Indeed, deleting an empty bag separates the decomposition tree into
components whose unions of bags are disjoint and have no graph edges
between them. Connectedness of $G$ implies that exactly one of these
components contains graph vertices, and all the others can be discarded.
Repeating this operation removes all empty bags.

\subsection{Separators and consistent bag laws}

\begin{lemma}[Separator properties]\label{lem:separator}
The following statements hold for a tree decomposition of a connected graph.
\begin{enumerate}
 \item If a bag $B$ lies on the tree path between bags $A$ and $C$, every
 graph path from a vertex of $A$ to a vertex of $C$ meets $B$.
 \item For adjacent bags $A,C$, their intersection separates the unions
 of bags on the two sides of the corresponding tree edge. In particular,
 adjacent nonempty bags have nonempty intersection.
 \item If $A_0,\ldots,A_r$ occur in this order on a tree path and
 $x\in A_r$, then $d(x,A_i)$ is nonincreasing in $i$.
 Moreover, if $B$ lies between $A$ and a bag containing $x$, some
 $v\in B$ on a shortest path from $A$ to $x$ satisfies
 \begin{equation}\label{eq:separator-geodesic}
  d(x,A)-d(v,A)=d(x,v).
 \end{equation}
\end{enumerate}
\end{lemma}

\begin{proof}
Delete the node corresponding to $B$ from $T$.
The occurrence subtree of each vertex outside $B$ lies in one resulting
component. Two adjacent graph vertices outside $B$ lie in the same
component, since some bag contains them both.
A graph path avoiding $B$ therefore stays in one component.
This proves the first assertion; if $B$ is an endpoint bag, that assertion
is immediate.

For a tree edge with endpoint bags $A,C$, a graph vertex occurs on both
sides precisely when it belongs to $A\cap C$, by connectedness of its
occurrence subtree. An edge joining the two sides away from this
intersection would have no bag containing both endpoints.
This proves the second assertion. If $A\cap C$ were empty, the nonempty
bags $A$ and $C$ would lie in different nonempty components of $G$.

For the third assertion, a shortest path from $A_i$ to $x$ meets
$A_{i+1}$ by the first assertion, so
$d(x,A_{i+1})\leq d(x,A_i)$.
Finally, choose a shortest path from $A$ to $x$ and let $v$ be any of its
vertices in $B$. Its initial segment has length $d(v,A)$: a shorter path
from $A$ to $v$ would shorten the chosen path to $x$.
The remaining segment has length $d(x,v)$, giving
\eqref{eq:separator-geodesic}.
\end{proof}

For a finite set $S$, write $\Omega_S=\{0,1\}^S$.
A \emph{labeled cut measure} on $S$ is a finite nonnegative measure on
$\Omega_S$. For $R\subseteq S$, its \emph{trace} $\nu|_R$ is the
pushforward under restriction to $R$.
The global complement map is $\omega\mapsto 1-\omega$; write $\nu^c$ for
its pushforward and
\[
 \nu^{\mathrm{sym}}=\frac{\nu+\nu^c}{2}
\]
for symmetrization. A measure is \emph{symmetric} if $\nu=\nu^c$.
The cut pseudometric represented by $\nu$ is
\begin{equation}\label{eq:cut-metric}
 D_\nu(u,v)=\nu\{\omega:\omega(u)\ne\omega(v)\}.
\end{equation}
Symmetrization leaves this pseudometric unchanged.

A family $\mu=(\mu_t)_{t\in V(T)}$ of bag measures is
\emph{consistent} if all its measures have the same mass $m>0$ and
\[
 \mu_t|_{B_t\cap B_s}=\mu_s|_{B_t\cap B_s}
 \qquad(ts\in E(T)).
\]
The common trace on an intersection $S=B_t\cap B_s$ is denoted by
$\mu_S$.

Sampling each bag conditionally on its overlap with the preceding bag
is the same Markov sampling principle used by Chlamt\'a\v{c},
Krauthgamer, and Raghavendra in their rounding procedure
\cite[Section~3.2, Lemmas~3.3--3.4]{CKR}.  We record its finite-measure
form, including the precise marginals needed below.

\begin{lemma}[Markov gluing]\label{lem:gluing}
Every consistent family of bag measures has a canonical Markov gluing
$\rho_\mu$ on $\Omega_V$, of mass $m$, with trace $\mu_t$ on each bag.
It is independent of the chosen root and is symmetric when every bag
measure is symmetric.
Along a path of bags, the successive intersection assignments, with a
vertex label from the corresponding endpoint bag at either end, form
a Markov chain under $\rho_\mu/m$.
The associated pseudometric
\[
 D_\mu(u,v)=\rho_\mu\{\omega:\omega(u)\ne\omega(v)\}
\]
is realized in a finite-dimensional real $L_1$ space.
\end{lemma}

\begin{proof}
Root $T$ and sample the root assignment with law $\mu_t/m$.
At a child bag $B_s$ with parent bag $B_t$, keep the labels on
$S=B_t\cap B_s$ and sample the remaining labels according to the
conditional law of $\mu_s$ given that assignment on $S$.
Previously assigned vertices in $B_s$ are exactly those in $S$: any
earlier occurrence of such a vertex is connected to $s$ through $t$.
Consequently the sampling defines an assignment on all of $V$.
An induction over the rooted tree, using consistency on $S$, proves that
each bag has marginal $\mu_s/m$.
Multiply the resulting probability measure by $m$ to obtain $\rho_\mu$.

For every full assignment $\omega$ for which all separator factors are
positive, its mass is
\begin{equation}\label{eq:gluing-product}
 \rho_\mu(\omega)
 =\frac{\displaystyle\prod_{t\in V(T)}\mu_t(\omega|_{B_t})}
        {\displaystyle\prod_{ts\in E(T)}
          \mu_{B_t\cap B_s}(\omega|_{B_t\cap B_s})}.
\end{equation}
If a separator assignment has zero mass, the full assignment has zero
mass as well, by the already established marginals.
Formula~\eqref{eq:gluing-product} is independent of the root and
invariant under complement when the bag measures are symmetric.

Rooting at the first bag of a path and integrating out branches shows
that each next intersection is sampled using only the previous
intersection assignment. This also permits an endpoint vertex as the
first or last layer, by conditioning the endpoint bag law on that
vertex or restricting it to that vertex, respectively.
Finally, on the finite measure space $(\Omega_V,\rho_\mu)$, the functions
$F_v(\omega)=\omega(v)$ satisfy
$\|F_u-F_v\|_1=D_\mu(u,v)$.
\end{proof}

Two immediate consequences will be used repeatedly.
First, if the separation mass in a bag containing each graph edge
$uv$ is at most $U\ell(uv)$, then the bag marginals and the triangle
inequality give
\begin{equation}\label{eq:edge-extension}
 D_\mu(u,v)\leq U d(u,v)\qquad(u,v\in V).
\end{equation}
Second, when the prepared laws depend measurably on auxiliary randomness
and have an integrable common mass, the average of $\rho_\mu$ is again
a finite cut measure. Thus averaging the resulting pseudometrics
preserves their $L_1$ representation.

\subsection{A transport inequality for a common submeasure}

The glued measure need not preserve the contribution of any one
function outside a bag.  The next result gives a substitute: if the
same cut measure remains present in every bag along a path segment,
its uncertainty at the terminal vertex forces separation in the glued
metric.  The published flow theorem supplies the bound independent of
the path's length; we prove the passage from a common submeasure to a
feasible flow.

The external input in this subsection is a theorem of
Chlamt\'a\v{c}, Krauthgamer, and Raghavendra.
A \emph{Markov flow graph} for a finite Markov chain
$X_0,\ldots,X_r$, with state sets $L_0,\ldots,L_r$, has arcs only from
$L_{i-1}$ to $L_i$, with capacities
\[
 c_i(a,b)=\Pr[X_{i-1}=a,\ X_i=b].
\]
We use the symmetric case: each layer has a fixed-point-free involution
$\theta_i$, and the joint law of $(\theta_i(X_i))_{i=0}^r$ equals that
of $(X_i)_{i=0}^r$.

\begin{theorem}[Symmetric Markov flow bound {
\cite[Theorem~4.4]{CKR}}]\label{thm:markov-flow}
For every integer $w\geq2$ there is a finite constant $C(w)\geq1$
with the following property.
Suppose a symmetric Markov flow graph has at most $w$ states per layer,
first layer $\{s_0,s_1\}$, and last layer $\{t_0,t_1\}$, and its
involutions interchange the two states in each endpoint layer.
Every capacity-respecting flow from $s_0$ to $t_1$ has value at most
\[
 C(w)\Pr[X_0=s_0,\ X_r=t_1].
\]
The constant is independent of the number of layers and of all
transition probabilities.
\end{theorem}

Zero-probability states may be removed in complementary pairs: their
incident capacities vanish, so neither feasible flows nor endpoint
probabilities change. Deterministic transitions are allowed.

For the rest of the proof fix
\begin{equation}\label{eq:bk-definition}
 b_k=C(2^k)^{-1}>0.
\end{equation}
For a labeled cut measure $\nu$ on $V$, a set $P\subseteq V$, and
$z\in V$, define
its \emph{uncertainty at $z$ given $P$} by
\begin{equation}\label{eq:uncertainty-definition}
 E_\nu(z\mid P)
 =\sum_{b\in\Omega_P}
   \min\bigl\{\nu(\omega|_P=b,\omega(z)=0),
              \nu(\omega|_P=b,\omega(z)=1)\bigr\}.
\end{equation}
Thus $E_\nu(z\mid P)$ is the total mass of the less frequent label at
$z$, summed over the possible assignments on $P$.  It measures the part
of the label at $z$ that cannot be predicted from that assignment.
This functional is nondecreasing in $\nu$, positively homogeneous, and
concave. Complementing all labels leaves it unchanged, so
\begin{equation}\label{eq:uncertainty-symmetrization}
 E_{\nu^{\mathrm{sym}}}(z\mid P)\geq E_\nu(z\mid P).
\end{equation}
These properties follow directly from the corresponding properties of
the minimum of two nonnegative coordinates; for complement invariance,
reindex the sum by $b\mapsto1-b$.

\begin{proposition}[Transport of uncertainty]\label{prop:uncertainty}
Let $\mu$ be a consistent family of symmetric bag measures, with bags of
size at most $k$. Consider the path $B_0,\ldots,B_r$ from a bag
containing $y$ to a bag containing $z$, and let $P=B_j$ lie on this path.
Suppose a labeled cut measure $\nu$ on $V$ satisfies
\[
 \nu|_{B_i}\leq\mu_{B_i}\qquad(j\leq i\leq r).
\]
Then the Markov-glued pseudometric satisfies
\begin{equation}\label{eq:uncertainty-transport}
 D_\mu(y,z)\geq b_k E_\nu(z\mid P).
\end{equation}
The measure $\nu$ need not be symmetric.
\end{proposition}

\begin{proof}
We will attach a prefix sampled from the glued law to a suffix sampled
from $\nu$.  It is enough to dominate each consecutive joint marginal;
no domination of $\nu$ by the full glued measure is asserted.  The full
assignment on $P$ is inserted as a layer so that the attachment uses
exactly the information in the uncertainty functional.

Write $\rho=\rho_\mu$ and let $m$ be its mass.
If $z\in P$, each assignment on $P$ determines the label at $z$, so
$E_\nu(z\mid P)=0$ and the assertion is immediate.
Otherwise, use the vertex sets
\[
 \{y\},\ B_0\cap B_1,\ldots,B_{r-1}\cap B_r,\ \{z\},
\]
as indices for layers: the states in the layer indexed by a set $R$
are the assignments in $\Omega_R$. Insert one additional layer
consisting of the full assignment on $P$. This layer goes between the incoming and outgoing intersections
at $P$; when $P=B_0$, it goes immediately after the $y$ layer.
The vertex sets indexing consecutive layers are contained in a single bag.

Under $\rho/m$ these layers form a Markov chain.
Indeed, the additional layer samples the rest of $P$ conditional on its
incoming assignment, and its outgoing intersection is then a
restriction of the full $P$ assignment. Summing over the additional
coordinates gives exactly the original intersection transition.
Every layer has at most $2^k$ states.
All indexing vertex sets are nonempty by Lemma~\ref{lem:separator}, so
complementation acts on their assignment spaces without fixed points.
Symmetry of the bag laws preserves the chain law.
The mass-$m$ capacity on each consecutive pair is the corresponding
projection of a bag measure.

We construct a measure $\lambda$ on paths in these layers.
Take the suffix measure, beginning with the full $P$ layer, to be the
projection of $\nu$. For its assignment $b$ on $P$, attach an
independent prefix with the conditional probability law under
$\rho/m$ given that same full assignment.
Put
\[
 m_b=\rho(\omega|_P=b),\qquad n_b=\nu(\omega|_P=b).
\]
The assumed domination at $P$ gives $0\leq n_b\leq m_b$.
For any prefix event $F$,
\[
 \lambda(F)
 =\sum_{b:m_b>0}n_b\Pr_\rho[F\mid\omega|_P=b]
 \leq\sum_{b:m_b>0}m_b\Pr_\rho[F\mid\omega|_P=b]
 =\rho(F).
\]
Here and below conditional probabilities under $\rho$ mean those under
$\rho/m$. Assignments with $m_b=0$ have $n_b=0$ and require no
conditional law.
Thus every prefix edge respects its original capacity.
On the suffix, the law is the corresponding projection of $\nu$;
because the vertex sets indexing consecutive layers lie within a single suffix bag, its
edges also respect the capacities.
The attachment is consistent on all overlapping vertex labels: a
vertex occurring on both sides of $P$ must occur in $P$, by the
connectedness of its occurrence subtree.

For $i\in\{0,1\}$ and $b\in\Omega_P$, set
\[
 a_i(b)=\nu(\omega|_P=b,\omega(z)=i),
 \qquad
 q_b=\Pr_\rho[\omega(y)=0\mid\omega|_P=b].
\]
For $m_b=0$, set $q_b=0$; both $a_i(b)$ then vanish.
The construction gives
\begin{align*}
 \lambda(\omega(y)\ne\omega(z))
 &=\sum_b\bigl(q_ba_1(b)+(1-q_b)a_0(b)\bigr)\\
 &\geq\sum_b\min\{a_0(b),a_1(b)\}
 =E_\nu(z\mid P).
\end{align*}
Restrict $\lambda$ first to paths with endpoint labels $(0,1)$, and
then to paths with endpoint labels $(1,0)$.
Each restriction defines a feasible flow: its arc masses satisfy
conservation at intermediate states and are dominated by the original
arc capacities.
Apply Theorem~\ref{thm:markov-flow} to the first flow and to the
complemented second flow, normalizing all masses by $m$.
After multiplying back by $m$ and adding, we obtain
\[
 \lambda(\omega(y)\ne\omega(z))
 \leq C(2^k)\rho(\omega(y)\ne\omega(z))
 =C(2^k)D_\mu(y,z).
\]
Together with the preceding lower bound and
\eqref{eq:bk-definition}, this proves the Proposition.
\end{proof}

\subsection{Threshold sweeps and height tests}

We now turn real functions into the common submeasures required by
Proposition~\ref{prop:uncertainty}.  The two tests below use either a
gap between function values or a pair of opposite perturbations that
agree on the testing set.

Fix a cutoff $M>0$.
For a function $F:V\to[-M,M]$, let $\sigma_F^M$ be the labeled cut
measure obtained by pushing Lebesgue measure on $[-M,M]$ through
\[
 t\longmapsto\omega_t,
 \qquad \omega_t(v)=\mathbf{1}_{\{F(v)>t\}}.
\]
Its total mass is $2M$. Its symmetric version is
$(\sigma_F^M)^{\mathrm{sym}}$, and both versions have separation mass
\begin{equation}\label{eq:sweep-separation}
 D_{\sigma_F^M}(u,v)=|F(u)-F(v)|.
\end{equation}
We say that $F$ has \emph{sweep margin} at least $a>0$ if
$|F(v)|\leq M-a$ for every $v\in V$.
The constant assignments produced near either cutoff are retained as
part of the measure.

\begin{lemma}[Height-gap test]\label{lem:height-gap}
Let $P\subseteq V$, $z\in V$, and $t>0$.
Suppose that $F(z)$ has distance at least $t$ from both endpoints of
$[-M,M]$ and that
\[
 |F(z)-F(s)|\geq t\qquad(s\in P).
\]
For $\nu=(\sigma_F^M)^{\mathrm{sym}}$, one has
\[
 E_\nu(z\mid P)\geq t.
\]
\end{lemma}

\begin{proof}
No value $F(s)$ with $s\in P$ lies in the open interval
$(F(z)-t,F(z)+t)$.
The assignment on $P$ is therefore constant throughout this threshold
interval. Its two halves each have length $t$, and the label at $z$
is opposite on the two halves. For that fixed assignment on $P$, the
unsymmetrized sweep gives mass at least $t$ to each label at $z$.
Thus $E_{\sigma_F^M}(z\mid P)\geq t$, and
\eqref{eq:uncertainty-symmetrization} proves the result.
\end{proof}

\begin{lemma}[Uncertainty at a paired launch]\label{lem:launch-gap}
Let $A\subseteq V$ and let $f,g:V\to\mathbb R$ satisfy $g|_A=0$ and
$f+g,f-g\in[-M,M]^V$.
Set
\[
 \nu=(\sigma_{f+g}^M)^{\mathrm{sym}}
       +(\sigma_{f-g}^M)^{\mathrm{sym}}.
\]
Then, for every $z\in V$,
\begin{equation}\label{eq:launch-uncertainty}
 E_\nu(z\mid A)\geq2|g(z)|.
\end{equation}
Moreover, for every $u,v\in V$,
\begin{equation}\label{eq:paired-separation}
 \begin{split}
 D_\nu(u,v)
 &=|f(u)-f(v)+g(u)-g(v)|\\
 &\quad+|f(u)-f(v)-g(u)+g(v)|
 \geq2|g(u)-g(v)|.
 \end{split}
\end{equation}
\end{lemma}

\begin{proof}
The two height functions agree on $A$.
For each threshold strictly between $f(z)-g(z)$ and $f(z)+g(z)$,
their assignments on $A$ consequently agree and their labels at $z$
are opposite.
For each $b\in\Omega_A$, let $I_b$ be the set of thresholds in this
interval giving assignment $b$ on $A$.
In the sum of the two unsymmetrized sweeps, each of the two possible
labels at $z$, together with assignment $b$ on $A$, has mass at least
the Lebesgue measure of $I_b$.
Summing over $b$ gives uncertainty at least the length of the whole
interval, namely $2|g(z)|$.
The sum of the symmetric sweeps is the symmetrization of their sum,
so \eqref{eq:uncertainty-symmetrization} proves
\eqref{eq:launch-uncertainty}.
Formula~\eqref{eq:sweep-separation} proves the equality in
\eqref{eq:paired-separation}; the final inequality follows from
$|a+b|+|a-b|\geq2|b|$.
\end{proof}

The two tests have the following consequences for the glued metric.
Suppose $g|_A=0$ and the symmetric sweep measures of $f+g$ and $f-g$
have a sum $\nu$ satisfying $\nu|_B\leq\mu_B$ at every bag on a path
segment from a bag $A$ to a bag containing $z$. If $A$ lies on the path
from a bag containing $y$ to that terminal bag, then
Proposition~\ref{prop:uncertainty} and Lemma~\ref{lem:launch-gap} give
\begin{equation}\label{eq:launch-transport}
 D_\mu(y,z)\geq2b_k|g(z)|.
\end{equation}
In particular, $|g(z)|\geq d(z,A)/2$ gives
$D_\mu(y,z)\geq b_kd(z,A)$.

Likewise, suppose a single function's symmetric sweep has trace dominated
by the bag laws on the segment beginning at a testing bag $P$ and ending
at a bag containing $z$, where $P$ lies on the path from a bag containing
$y$ to that terminal bag. If auxiliary randomness produces a height gap
$t$ as in Lemma~\ref{lem:height-gap} with probability at least $p$,
and the trace domination holds on that event, apply the two estimates to
each realization in the event. Averaging the pointwise inequality gives
\begin{equation}\label{eq:averaged-height-test}
 \mathbb E D_\mu(y,z)\geq b_kpt.
\end{equation}
The gap event may depend on the function values; the requirement
is that the tested sweep remain a common submeasure throughout the
specified segment on that event. Section~\ref{sec:particles} constructs
bag laws with precisely this persistence property.
